% MỞ ĐẦU

\chapter*{MỞ ĐẦU}
\addcontentsline{toc}{chapter}{MỞ ĐẦU}

\label{Chapter0}

%----------------------------------------------------------------------------------------

U não là một trong những loại ung thư nguy hiểm nhất với tỷ lệ sống sót sau năm năm thuộc hàng thấp nhất trong nhóm các bệnh ung thư. Chẩn đoán sớm và chính xác đóng vai trò quyết định trong việc lựa chọn phác đồ điều trị và cải thiện tình trạng cho bệnh nhân. Ảnh cộng hưởng từ (MRI) là phương tiện chẩn đoán hình ảnh chủ đạo cho bệnh lý não bộ, song việc phân tích thủ công các lát cắt MRI đòi hỏi nhiều thời gian và phụ thuộc vào kinh nghiệm của bác sĩ chuyên khoa.

Báo cáo này trình bày thiết kế và triển khai hệ thống BrainSeg-XAI -- một hệ thống trí tuệ nhân tạo hỗ trợ chẩn đoán u não từ ảnh MRI, kết hợp khả năng phân đoạn vùng khối u tự động với cơ chế giải thích kết quả dự đoán có thể hiểu được. Hệ thống được xây dựng dựa trên kiến trúc học sâu U-Net và U-Net++, tích hợp kỹ thuật Gradient-weighted Class Activation Mapping (Grad-CAM) để trực quan hóa lý do mô hình đưa ra quyết định, và một bộ máy luật y khoa để đánh giá mức độ rủi ro lâm sàng.

Đề tài thuộc lĩnh vực giao thoa giữa học máy, thị giác máy tính và tin học y tế -- ba lĩnh vực đang có sự phát triển mạnh mẽ trong bối cảnh ứng dụng trí tuệ nhân tạo vào thực tiễn lâm sàng.
